\section{Technology-induced screening and the value of public access}\label{sec:diagnostic}

\subsection{Central proposition: adverse local innovation and valuable access}
Consider the public-access margin from private-only provision ($P$) to free basic access coexisting with private premium provision ($B$). Fix the partner jurisdiction in $G$ or $B$, so $(m_j,q_j)=(1,1/2)$; the private intermediary in region $i$ nevertheless reoptimises its fee in the private-only continuation. This is a \emph{conditional institutional comparison}, not a global comparative static of all four stage-1 entry decisions.

Write $d=\Delta_i$, $\gamma=\rho\mu$, $C=\gamma N_iN_j$, and set
\begin{equation}
 U=u_i+\gamma N_j\omega_iR_0,\quad
 \eta_i=\frac{\gamma N_j\omega_i\Pi}{2},\quad
 D=d+\eta_i\lambda,\quad
 \ell=\frac{D-U}{2D},
 \label{eq:screen-primitive}
\end{equation}
on the interior branch $0<U<D$. The private-only fee is $(U+D)/2$, whereas free basic provision restores access to all founder types, with premium pricing at $d/2$.
Using Lemma~3, the gross local public-access value, before government fixed cost, is
\begin{equation}
 T_i(\lambda)=N_i\left(u_i\ell-\frac d8+\frac{d\ell^2}{2}\right)
        +C\left(S_i^0\ell+\frac{\lambda K_i\ell^2}{4}\right),
 \qquad K_i=M_i+\omega_i\Pi .
 \label{eq:screen-threshold}
\end{equation}
Differentiating while keeping the research-market primitives, $\rho$, and $u_i$ fixed gives the exact identity
\begin{equation}
 \begin{split}
 T_i'(\lambda)={}&
 \left[N_i(u_i+d\ell)+CS_i^0\right]
       \frac{U\eta_i}{2D^2}\\
 &+\frac{CK_i}{4}
       \left(\ell^2+\frac{\lambda\ell U\eta_i}{D^2}\right).
 \end{split}
 \label{eq:diagnostic-derivative}
\end{equation}
The first line describes the access restored by the rising private fee-induced exclusion cutoff, including its baseline matching component; the second is the regional return to the technically different additional research pairs. The two need not have the same sign.

\paragraph{Proposition 1 (the screening--incidence inversion).}\label{prop:inversion}
Maintain the model's uniform and perfectly aligned technical-skill/premium-value distribution, a positive interior Cournot continuation, an interior local private-only fee, partner full participation, $M_i<0$, $0<u_i<d$, $\Pi>0$, and nonnegative $R_0$. Fix a compact interval of technologically feasible $\lambda\geq0$. Define
\begin{equation}
 \omega_i^{\mathrm{crit}}
 =\frac{-M_i}{\Pi}
     \frac{(d-u_i)^2}{d^2+3u_i^2}.
 \label{eq:omega-critical}
\end{equation}
Suppose the interval
\begin{equation}
 \boxed{\qquad
 \omega_i^{\mathrm{crit}}<\omega_i<
            \min\left\{1,\frac{-M_i}{\Pi}\right\}
 \qquad}
 \label{eq:screen-inversion-band}
\end{equation}
is nonempty. As $\gamma=\rho\mu\downarrow0$, holding $\rho$ and the basic-service value $u_i$ fixed, the conditional public-access threshold satisfies, uniformly on the specified $\lambda$ interval,
\begin{equation}
 \frac{\partial T_i}{\partial\lambda}
  =\frac{\gamma N_iN_j}{16d^2}
    \left[M_i(d-u_i)^2+
          \omega_i\Pi(d^2+3u_i^2)\right]
       +O(\gamma^2).
 \label{eq:main-primitive-slope}
\end{equation}
Consequently, at every sufficiently small positive $\gamma$,
\begin{equation}
 K_i<0,\qquad S_i^0<0,\qquad
 \frac{\partial T_i}{\partial\lambda}>0
 \label{eq:main-inversion}
\end{equation}
throughout the $\lambda$ interval. If in addition $0<u_i<2d/3$, then $T_i(\lambda)>0$ uniformly at sufficiently small positive $\gamma$. Thus public provision can have a positive affordable fixed cost and become more valuable as technical complementarity rises even though its incremental research-commercialisation surplus is strictly negative in the financing region.

The critical share $\omega_i^{\mathrm{crit}}$ strictly decreases in $u_i/d$ on $(0,1)$. At $\omega_i=0$ and $M_i<0$, by contrast, the private screening response to $\lambda$ disappears and $T_i'(\lambda)=CK_i\ell^2/4<0$ on the interior branch. These are implications of the existing primitives rather than new policy instruments. The proof and the uniformity conditions are in the Online Supplement. The proposition concerns the \emph{gross conditional public cost ceiling}, not an all-parameter four-player subgame-perfect equilibrium or an optimum across taxes, subsidies and physical facilities.

Two distinctions matter. First, $S_i^0=p_bK_i-\omega_i\zeta^2\Pi^2/(2k)<0$ when $K_i<0$; the adverse regional return applies to baseline as well as skill-complementarity-induced ventures. Second, \eqref{eq:screen-inversion-band} provides a nonempty feasible band only when $\omega_i^{\mathrm{crit}}<\min\{1,-M_i/\Pi\}$; positive basic-service value alone does not guarantee that its intersection with feasible profit shares is nonempty. The sign restriction is first-order in weak matching, while Section~\ref{sec:numerical} separately certifies a finite-matching open rectangle for the original economic environment.

\subsection{Equal-mean homogeneous success and repriced founder selection}
A related diagnostic explains why matching technologies with equal full-population mean success rates can imply different values of public access. Compare the heterogeneous premium $\lambda\theta_i\theta_j$ ($H$) with a pooled homogeneous premium $\lambda/4$ ($L$). With the same partner fully participating, write $h=\eta_i\lambda$ and retain $U$ from \eqref{eq:screen-primitive}. Their respective globally optimal private-only cutoffs, including the full-coverage corners, are
\begin{equation}
 \ell_H=\max\left\{0,\frac{d+h-U}{2(d+h)}\right\},\qquad
 \ell_L=\max\left\{0,\frac{d-U-h/2}{2d}\right\}.
 \label{eq:diagnostic-cuts}
\end{equation}
The pooled technology increases the private-only participation intercept by $h/2$ and removes the slope $h$; the two comparisons must not be evaluated at the same unoptimised fee. For $h>0$, $\ell_H\geq\ell_L$, with strict inequality unless both cutoffs are zero. When both are interior,
\begin{equation}
 \ell_H-\ell_L=\frac{h(d+h+2U)}{4d(d+h)}>0.
 \label{eq:diagnostic-cutgap}
\end{equation}

Let $T_H$ and $T_L$ denote the public cost ceilings for the heterogeneous and equal-mean pooled technologies after optimal fee readjustment. With $\delta=\ell_H-\ell_L$, substitution into Lemma~3 gives
\begin{equation}
 \begin{split}
 T_H-T_L={}&\delta\left[N_i\left(u_i+\frac d2(\ell_H+\ell_L)\right)
                         +CS_i^0\right]\\
 &+\frac{C\lambda K_i}{4}(\ell_H^2-\ell_L).
 \end{split}
 \label{eq:diagnostic-gap}
\end{equation}
This holds also at the full-participation corners. The first line measures restored basic-service value and baseline matched-project incidence when privately selected founder exclusion changes. The second compares different skill-weighted exposure of restored projects. Its sign can differ from the first line, so even equal full-participation mean project success need not imply equal equilibrium access value. If a fixed government cost lies strictly between $T_H$ and $T_L$, the conditional public-entry decision reverses; this alone does not certify the global four-entrant equilibrium.

The comparison can also reverse in the other direction. With $\omega_i=0$, $M_i>0$ and $0<u_i<d$, the two technologies induce the same positive cutoff $\ell$, but
$T_H-T_L=C\lambda M_i\ell(\ell-1)/4<0$ for $\lambda>0$.
The strict sign persists for sufficiently small positive $\omega_i$ by continuity. A positive technical effect on local product-market welfare can therefore diminish the marginal value of restoring a lower-skill group. The proposition and this equal-mean diagnostic answer different questions: Proposition~1 changes $\lambda$ within the heterogeneous technology; \eqref{eq:diagnostic-gap} changes the distribution of technical gains at fixed $\lambda$. Neither yields an unrestricted policy prescription.